\chapter{Cross-Validation against GMAT and Orekit}
\label{chap:cross-validation}

Every model specified in Part~\ref{part:math} is checked against at least one
independent implementation. Three references are used: Orekit v13.1, a
JVM-based astrodynamics library \citep{orekit2025}; GMAT R2026a,
NASA's General Mission Analysis Tool \citep{gmat2023}, pointed at \lupnt{}'s own
DE440 ephemeris (see \cref{sec:xval-constants}); and NAIF's CSPICE
toolkit \citep{acton1996spice}, which \lupnt{} links directly and can
therefore query in-process. The first two are compared through checked-in JSON
fixtures; the third is compared live, inside the ordinary test run.

This chapter states how each comparison is constructed, tabulates every
measured difference and the tolerance asserted against it, and --- the part
that matters --- explains at specification level \emph{why} each residual is
the size it is. A tolerance that merely passes is not the goal. The intent is
that no residual in the tables below is unexplained: each one is attributed
either to floating-point representation, to a difference in the input data
(EOP vintage, DE version), or to a deliberate modelling choice
(inertial-axis $J_2$, mean-Earth versus principal-axes lunar frame). Where a
difference traces back to a model, the chapter that specifies that model is
cited: \cref{chap:time-conversions} for the time scales,
\cref{chap:frame-conversions} for the rotations, and \cref{chap:dynamics} for
the force models.

The implementing files are \file{cpp/test/orekit/} and \file{cpp/test/gmat/}
(fixture-based), and \file{cpp/test/conversions/test_frame_conversions.cc} and
\file{cpp/test/conversions/test_time_conversions_ttdb.cc} (live SPICE). The
external data products the comparison depends on are specified in
\cref{chap:sp3-download}.

% ===========================================================================
\section{Methodology}
\label{sec:xval-method}
% ===========================================================================

Neither Orekit nor GMAT is a build, run, or test dependency of \lupnt{}.
A two-stage arrangement keeps it that way:

\begin{enumerate}[leftmargin=1.6em]
  \item A developer-only \emph{generator} drives the external tool once and
    serialises its outputs into a checked-in JSON fixture.
    \file{cpp/test/orekit/gen_orekit_reference.py}, run in the optional
    \code{dev} pixi environment, writes
    \file{cpp/test/orekit/data/orekit_reference.json}.
    \file{cpp/test/gmat/gen_gmat_reference.py}, which needs a local GMAT
    desktop installation, generates a GMAT script, runs it, parses the report
    output, and writes \file{cpp/test/gmat/data/gmat_reference.json}. The
    exact script it ran is checked in beside the fixture as
    \file{cpp/test/gmat/data/gmat_reference.script}, with machine-specific
    paths replaced by placeholder tokens.
  \item Ordinary Catch2 test files under \file{cpp/test/orekit/} and
    \file{cpp/test/gmat/} load the fixture through \cpp{LoadTestJson()} in
    \file{cpp/test/utils.cc} and compare \lupnt{}'s own computation against
    it. They run as part of \code{pixi run test-cpp} on every machine, with no
    JVM, no GMAT, and no network.
\end{enumerate}

\subsection{Isolating algorithmic differences}
\label{sec:xval-constants}

Both fixtures are generated with \lupnt{}'s physical constants forced onto the
external tool, so a residual can never be a constants mismatch in disguise.
The values, from \file{cpp/lupnt/core/constants.h}, are listed in
\cref{tab:xval-constants} and recorded in the \code{\_meta} block of each
fixture.

\begin{table}[htbp]
  \centering
  \caption{Constants imposed on both sides of every comparison.}
  \label{tab:xval-constants}
  \small
  \begin{tabular}{@{}llll@{}}
    \toprule
    Symbol & \lupnt{} identifier & Value & Unit \\
    \midrule
    $\mu_\oplus$ & \code{GM\_EARTH}  & \num{398600.435507e9} & \si{\meter\cubed\per\second\squared} \\
    $\mu_{\mathrm{M}}$& \code{GM\_MOON}   & \num{4902.800118e9}   & \si{\meter\cubed\per\second\squared} \\
    $R_\oplus$   & \code{R\_EARTH}   & \num{6378137}         & \si{\meter} \\
    $J_2$        & \code{J2\_EARTH}  & \num{1.08262668e-3}   & --- \\
    \bottomrule
  \end{tabular}
\end{table}

On the GMAT side this required more than a configuration flag. A minimal
potential file, \file{cpp/test/gmat/lupnt_j2_earth.cof}, sets \code{GM},
\code{R\_eq} and the normalised \code{C20} $(= -J_2$ normalised$)$ to exactly
these values; \code{TwoBodyFM} uses it at \code{Degree=0, Order=0} (pure point
mass) and \code{J2FM} at \code{Degree=2, Order=0} (point mass plus the
degree-2 zonal term). \code{Earth.Mu}, \code{Earth.EquatorialRadius} and
\code{Luna.Mu} are additionally overridden globally so that GMAT's internal
Cartesian--Keplerian conversions use the same numbers.

The planetary ephemeris is now shared as well. GMAT ships only up to DE424 ---
no DE440 even in R2026a --- so both generators point it at \lupnt{}'s own
\file{de440.bsp} through \code{SolarSystem.EphemerisSource = 'SPICE'} and
\code{SolarSystem.SPKFilename}. Orekit likewise evaluates DE440
\citep{park2021de440}. With the same ephemeris on every side, the former
DE421-versus-DE440 Moon-frame offset is gone: what remains in the Moon-centred
and third-body comparisons is a distribution and interpolation residual, not a
DE-version difference. For the third-body and SRP force models this also means
the Sun and Moon positions that drive the perturbation are identical on both
sides.

Where the two comparisons overlap --- the frame-conversion input states, the
orbit cases, the epochs --- the fixtures use \emph{identical} inputs. The
Orekit and GMAT columns in the tables below are therefore directly comparable,
and the difference between the two columns is itself informative: it is what
lets \cref{sec:xval-causes} attribute the $J_2$ propagation residual to the
symmetry-axis convention rather than to the integrator.

\subsection{Tolerance philosophy}
\label{sec:xval-tolerances}

Each asserted tolerance is set a small factor --- typically $2$ to $20$ ---
above the measured difference, with a comment in the test file naming the
origin of that difference. Tolerances that scale with the state are written
that way rather than flattened to a constant: the GCRF--ITRF position
tolerance is $\max(\SI{100}{\meter},\, \num{2e-5}\,\lVert\vec r\rVert)$
because the underlying cause is a small rotation angle, whose linear effect
grows with radius. Flattening it would make the LEO test vacuous and the
trans-lunar test unpassable.

% ===========================================================================
\section{Scope of the comparison}
\label{sec:xval-scope}
% ===========================================================================

\begin{table}[htbp]
  \centering
  \caption{What each external reference covers.}
  \label{tab:xval-scope}
  \small
  \begin{tabularx}{\textwidth}{@{}l L L@{}}
    \toprule
    Category & versus Orekit & versus GMAT \\
    \midrule
    Time scales &
      $\TAI-\UTC$, $\TT-\TAI$, $\TDB-\TAI$, $\GPST-\TAI$, $\UTone-\UTC$ at six
      epochs bracketing the 2012 and 2016 leap seconds &
      $\TAI-\UTC$, $\TT-\TAI$, $\TDB-\TAI$ at the same epochs (GMAT R2026a
      exposes no \code{UT1ModJulian} or \code{GPSModJulian} parameter) \\
    Sidereal angles &
      GMST (IAU-82 / IERS 1996) and the Earth Rotation Angle (IAU 2000) &
      --- \\
    Earth frames &
      GCRF $\leftrightarrow$ EME2000 and GCRF $\leftrightarrow$ ITRF (IERS
      2010, no \code{simpleEOP}) for LEO, MEO, GEO and trans-lunar states at
      three epochs &
      the same states and epochs through GMAT's \code{EarthICRFCS},
      \code{EarthMJ2000Eq} and \code{EarthFixedCS} systems \\
    Moon frames &
      GCRF $\leftrightarrow$ \code{MOON\_CI} (DE440 translation);
      \code{MOON\_ME} against the analytical IAU pole model &
      GCRF $\leftrightarrow$ \code{MOON\_CI} (DE440 via SPICE); \code{MOON\_PA}
      against GMAT's SPICE-kernel Luna \code{BodyFixed}, a principal-axes
      frame \\
    Ephemerides &
      Earth$\to$Moon and Earth$\to$Sun position and velocity (DE440) at three
      epochs &
      implicit in the Moon-frame translations \\
    Kepler propagation &
      six orbits --- Earth MEO-HEO, LEO, GEO, Molniya plus lunar LLO and ELFO
      --- propagated to $0.25$, $0.5$, $0.75$ and $1.5$ periods with the
      analytical \longid{KeplerianPropagator} &
      the same six orbits, numerically propagated point-mass (RK89,
      \code{Accuracy = 1e-13}) \\
    $J_2$ acceleration &
      formula-level check at five positions
      (\longid{J2OnlyPerturbation.computeAccelerationInJ2Frame}) &
      --- \\
    $J_2$ propagation &
      three orbits with $J_2$ about the \textbf{inertial} $Z$ axis, the same
      model as \lupnt{} --- isolates integrator differences &
      two orbits with $J_2$ evaluated in the \textbf{body-fixed} frame ---
      exposes the modelling difference \\
    Harmonic gravity &
      formula-level acceleration at 42 body-fixed positions (Earth $8\times8$
      EGM96, Moon $12\times12$ GRGM1200B) via
      \longid{HolmesFeatherstoneAttractionModel}, both sides fed the same
      \code{.cof} through an ICGEM bridge &
      the Earth $8\times8$ field propagated \SI{2}{\hour} with the shared
      \code{EGM96.cof} \\
    Third body &
      formula-level acceleration from the Sun and Moon
      (\longid{ThirdBodyAttraction}), perturber positions from the shared
      DE440 &
      Sun and Moon point masses propagated \SI{2}{\hour} \\
    Solar radiation pressure &
      formula-level cannonball acceleration, fully sunlit
      (\longid{SolarRadiationPressure}) &
      cannonball propagated \SI{2}{\hour} on a continuously sunlit orbit \\
    \bottomrule
  \end{tabularx}
\end{table}

The Orekit fixture is organised into the sections \code{time\_scales},
\code{sidereal}, \code{frames}, \code{moon\_frames}, \code{ephemerides},
\code{kepler}, \code{j2\_acceleration}, \code{j2\_propagation},
\code{gravity\_acceleration}, \code{third\_body\_acceleration} and
\code{srp\_acceleration}; the GMAT fixture
\file{cpp/test/gmat/data/gmat_reference.json} into \code{time\_scales},
\code{frames}, \code{moon\_frames}, \code{kepler} and \code{j2\_propagation}.
The GMAT force-model checks live in a separate fixture
\file{cpp/test/gmat/data/gmat_forces_reference.json} (sections
\code{gravity\_propagation}, \code{third\_body\_propagation} and
\code{srp\_propagation}), generated by
\file{cpp/test/gmat/gen_gmat_forces_reference.py} so the original
kepler/$J_2$/frame fixture is untouched. Within a \code{kepler} or
\code{j2\_propagation} entry, \code{coe0} is
$[a, e, i, \Omega, \omega, M_0]$ in SI units and radians with \emph{mean}
anomaly, \code{body} is \code{"EARTH"} or \code{"MOON"}, and \code{cart0} and
each \code{propagated[].cart} are $[x,\dots,v_z]$ in metres and metres per
second --- expressed in \code{EarthMJ2000Eq} axes for the Earth cases and
Luna-centred ICRF axes for the lunar ones.

\begin{lupntnote}
GMAT's Keplerian display state accepts only \emph{true} anomaly as its sixth
element, so \file{cpp/test/gmat/gen_gmat_reference.py} solves Kepler's
equation itself and uses the resulting true anomaly purely to configure each
spacecraft. The \code{coe0}, anomaly and \code{period} entries in the GMAT
fixture are therefore Python-computed; \code{cart0} and every
\code{propagated} entry are genuine GMAT outputs.
\end{lupntnote}

% ===========================================================================
\section{Measured differences}
\label{sec:xval-results}
% ===========================================================================

The numbers below were measured when the fixtures were generated (June 2026).
They are properties of a particular pair of data vintages, not constants of
nature: regenerate the fixtures (\cref{sec:xval-regenerate}) and re-derive
them whenever the comparison data changes. ``Tolerance'' is the value the
tests actually assert.

\subsection{Time scales and sidereal angles}
\label{sec:xval-time}

\begin{table}[htbp]
  \centering
  \caption{Time-scale and Earth-rotation-angle differences. Implemented in
    \code{cpp/test/orekit/test\_time\_conversions\_orekit.cc} and
    \code{cpp/test/gmat/test\_time\_conversions\_gmat.cc}; the models are
    specified in \cref{chap:time-conversions}.}
  \label{tab:xval-time}
  \small
  \begin{tabularx}{\textwidth}{@{}W{0.85} r r W{1.15}@{}}
    \toprule
    Quantity & versus Orekit (max) & versus GMAT (max) & Tolerance and origin \\
    \midrule
    $\TAI-\UTC$, $\TT-\TAI$, $\GPST-\TAI$ &
      ${\sim}\SI{2e-8}{\second}$ & ${\sim}\SI{6e-7}{\second}$ &
      \SI{1e-6}{\second}. Exactly defined offsets, so the residual is pure
      floating-point representation: \lupnt{} stores epochs as
      seconds-since-J2000 doubles, whereas GMAT reports ${\sim}30\,000$-day
      Modified Julian dates whose subtraction amplifies the ULP. \\
    $\TDB-\TAI$ &
      ${\sim}\SI{4e-8}{\second}$ & ${\sim}\SI{5e-7}{\second}$ &
      \SI{1e-6}{\second} (Orekit) / \SI{1e-5}{\second} (GMAT). The
      ${\sim}\SI{1.7}{\milli\second}$ periodic relativistic term is
      implemented with different truncated series on each side. \\
    $\UTone-\UTC$ &
      $\leq \SI{0.024}{\second}$ (away from leap seconds) & --- &
      \SI{0.1}{\second}. EOP-table vintage difference: the two sides hold
      different IERS bulletin snapshots. See \cref{sec:xval-leap}. \\
    GMST, ERA &
      $\leq \SI{1.8e-6}{\radian}$ & --- &
      \SI{1e-5}{\radian}. Identical formulas on both sides (IAU-82 GMST,
      IAU-2000 ERA); the residual is exactly the $\UTone$ difference times the
      Earth rotation rate. \\
    \bottomrule
  \end{tabularx}
\end{table}

\subsection{Reference frames and ephemerides}
\label{sec:xval-frames}

\begin{table}[htbp]
  \centering
  \caption{Frame-conversion and ephemeris differences. Implemented in
    \code{cpp/test/orekit/test\_frame\_conversions\_orekit.cc},
    \code{cpp/test/orekit/test\_ephemeris\_orekit.cc} and
    \code{cpp/test/gmat/test\_frame\_conversions\_gmat.cc}; the models are
    specified in \cref{chap:frame-conversions}. $\lVert\vec r\rVert$ is the
    position magnitude of the test state.}
  \label{tab:xval-frames}
  \scriptsize
  \begin{tabularx}{\textwidth}{@{}l L L L@{}}
    \toprule
    Conversion & versus Orekit (max) & versus GMAT (max) & Tolerance and origin \\
    \midrule
    GCRF $\leftrightarrow$ EME2000 &
      \SI{1.8e-4}{\meter} (trans-lunar), i.e.\
      ${\sim}\num{5e-13}\,\lVert\vec r\rVert$ &
      \SI{54}{\meter} (trans-lunar), i.e.\
      ${\sim}\num{1.6e-7}\,\lVert\vec r\rVert$ &
      \SI{1e-3}{\meter} (Orekit);
      $\max(\SI{5}{\meter},\, \num{3e-7}\,\lVert\vec r\rVert)$ (GMAT).
      \lupnt{} and Orekit implement the same IERS frame-bias matrix
      \citep{iers2010}; GMAT's ICRF $\leftrightarrow$ MJ2000Eq rotation
      differs by ${\sim}\SI{1e-7}{\radian}$, slightly epoch-dependent. \\
    GCRF $\leftrightarrow$ ITRF &
      \SI{4.4}{\kilo\meter} at the trans-lunar 2025 epoch
      (${\sim}\num{1.3e-5}\,\lVert\vec r\rVert$);
      \SIrange{14}{98}{\meter} at LEO &
      \SI{1.6}{\kilo\meter} trans-lunar
      (${\sim}\num{4.5e-6}\,\lVert\vec r\rVert$);
      \SIrange{34}{36}{\meter} at LEO &
      $\max(\SI{100}{\meter},\, \num{2e-5}\,\lVert\vec r\rVert)$, velocity
      \SI{1}{\meter\per\second}. Earth-orientation ($\UTone$, polar motion)
      table vintage: a tens-of-milliseconds $\UTone$ difference rotates
      positions by $\num{1e-6}$ to $\num{1e-5}$ radians about the polar axis,
      scaling with radius. Differences grow toward 2025 epochs, where one or
      both tables are predictions. \\
    GCRF $\rightarrow$ \code{MOON\_CI} &
      \SI{0.33}{\meter} & distribution level &
      \SI{1}{\meter} (Orekit) / \SI{200}{\meter} (GMAT). A pure DE-ephemeris
      translation: \lupnt{}, Orekit and now GMAT R2026a (through \lupnt{}'s
      \file{de440.bsp} via SPICE) all evaluate DE440 \citep{park2021de440}, so
      every column is a distribution and interpolation residual at
      ${\sim}\num{5e-11}$ relative, not a DE-version offset. The \SI{200}{\meter}
      GMAT tolerance predates the shared ephemeris and is retained with wide
      margin. \\
    Moon body-fixed &
      \code{MOON\_ME} versus the IAU model: \SIrange{57}{200}{\meter}
      (${\sim}\num{3e-5}\,\lVert\vec r\rVert$) &
      \code{MOON\_PA} versus Luna \code{BodyFixed}: distribution level,
      orientation-limited &
      $\num{3e-4}\,\lVert\vec r\rVert$ (Orekit) / \SI{300}{\meter} (GMAT).
      Orekit's body-oriented Moon frame is the analytical IAU pole model,
      which approximates the DE \emph{mean-Earth} axes to
      ${\sim}\SI{3e-5}{\radian}$. GMAT loads a SPICE Luna principal-axes
      kernel, which validates \lupnt{}'s \code{MOON\_PA}: with both sides on
      DE440 the Moon-position translation agrees at the distribution level and
      the PA orientations agree at the sub-arcsecond level, so the residual is
      well inside the retained \SI{300}{\meter} tolerance. Cross-check: swapping
      the frames (PA against IAU, or ME against GMAT) increases the differences
      by a factor of $10$ to $40$ (\SIrange{0.6}{3.4}{\kilo\meter}), confirming
      each identification. \\
    Ephemerides (Earth$\to$Moon, Earth$\to$Sun) &
      Moon \SI{0.33}{\meter}, Sun \SI{7.2}{\meter}, both
      ${\sim}\num{5e-11}$ relative & --- &
      $\max(\SI{1}{\meter},\, \num{1e-10}\,\lVert\vec r\rVert)$, velocity
      \SI{1e-5}{\meter\per\second}. The same DE440 data
      \citep{park2021de440}, independently distributed and independently
      interpolated. \\
    \bottomrule
  \end{tabularx}
\end{table}

\subsection{Orbit propagation}
\label{sec:xval-orbits}

\begin{table}[htbp]
  \centering
  \caption{Propagation differences. Implemented in
    \code{cpp/test/orekit/test\_orbit\_dynamics\_orekit.cc} (test cases
    \code{dynamics.kepler\_orekit\_reference},
    \code{dynamics.j2\_acceleration\_orekit\_reference},
    \code{dynamics.j2\_propagation\_orekit\_reference}) and
    \code{cpp/test/gmat/test\_orbit\_dynamics\_gmat.cc}
    (\code{dynamics.kepler\_gmat\_reference},
    \code{dynamics.j2\_propagation\_gmat\_reference}); the models are
    specified in \cref{chap:dynamics}.}
  \label{tab:xval-orbits}
  \scriptsize
  \begin{tabularx}{\textwidth}{@{}l L L L@{}}
    \toprule
    Quantity & versus Orekit (max) & versus GMAT (max) & Tolerance and origin \\
    \midrule
    Kepler: initial COE $\rightarrow$ Cartesian &
      \SI{1.5e-8}{\meter} & \SI{1.3e-8}{\meter} &
      \SI{1e-3}{\meter}. Pure conversion arithmetic; machine precision. \\
    Kepler: propagated Cartesian &
      \SI{5.7e-8}{\meter} / \SI{5e-12}{\meter\per\second} &
      \SI{1.8e-2}{\meter} / \SI{8e-6}{\meter\per\second} &
      \SI{1e-2}{\meter} (Orekit) / \SI{0.1}{\meter} (GMAT). Orekit's
      propagator is analytical, like \lupnt{}'s, so agreement is at machine
      precision; GMAT integrates even the point-mass problem numerically
      (RK89, accuracy \num{1e-13}), so its column is integrator drift. \\
    Kepler: propagated angles ($M$, $E$, $\nu$) &
      $\leq \SI{3e-15}{\radian}$ & $\leq \SI{5e-8}{\radian}$ &
      \SI{1e-9}{\radian} (Orekit) / \SI{1e-6}{\radian} (GMAT). The same
      analytical-versus-numerical distinction. \\
    $J_2$ acceleration &
      $\leq \SI{1e-15}{\meter\per\second\squared}$ &
      --- &
      \SI{1e-12}{\meter\per\second\squared}. Identical closed-form $J_2$
      expression on both sides. \\
    $J_2$ propagation, inertial-axis $J_2$ &
      \SI{2.2e-5}{\meter} / \SI{1.9e-9}{\meter\per\second} over 1.5 orbits &
      --- &
      \SI{1e-3}{\meter} / \SI{1e-6}{\meter\per\second}. The same force model
      on both sides ($J_2$ about the inertial $Z$ axis); the residual is
      \lupnt{}'s fixed-step RK4 at \SI{1}{\second} against Orekit's adaptive
      DP853. \\
    $J_2$ propagation, body-fixed $J_2$ &
      --- &
      \SIrange{22}{182}{\meter} /
      \SIrange{0.005}{0.099}{\meter\per\second} over $0.25$ to $1.5$ orbits &
      \SI{300}{\meter} / \SI{0.15}{\meter\per\second}. \textbf{Modelling
      difference}: \cpp{JToCartTwoBodyDynamics} takes the inertial $Z$ axis as
      the oblateness symmetry axis, whereas GMAT evaluates the zonal harmonic
      in the Earth body-fixed frame. The Orekit row above --- same dynamics,
      same axis convention, ${\sim}\SI{1e-5}{\meter}$ agreement --- isolates
      the axis choice as the cause. \\
    \bottomrule
  \end{tabularx}
\end{table}

\subsection{Force models}
\label{sec:xval-forces}

Three force models are cross-validated with two deliberately different
constructions, because the two reference tools expose different quantities.
Orekit exposes force-model \emph{accelerations}, so \lupnt{}'s acceleration
functions --- \cpp{AccelarationGravityField}, \cpp{AccelerationPointMass} and
\cpp{AccelerationSolarRadiation}, specified in \cref{chap:dynamics} --- are fed
the exact inputs Orekit's models used and the acceleration vector is compared
directly. There is no integrator or frame realisation in the loop: the check
isolates the force-model \emph{algorithm}, the same style as the
$J_2$-acceleration row of \cref{tab:xval-orbits}. GMAT reports only states, so
its checks are \emph{propagation}-level: a spacecraft is propagated for a fixed
\SI{2}{\hour} arc with one force model active and the Cartesian state is
compared, exercising the assembled propagation (force model plus integrator plus
frame). \Cref{tab:xval-forces} gives both.

For the Orekit gravity check both sides read the \emph{same} \lupnt{} \code{.cof}
coefficient file: the generator parses it and writes an Orekit-readable ICGEM
copy, so the comparison isolates the harmonic recursion and the normalisation
convention rather than the adopted coefficients. For third body and SRP the
perturber position (Sun, Moon) and the reference solar pressure are stored in
the fixture and fed to \lupnt{} unchanged, decoupling the ephemeris --- already
checked by the \code{ephemerides} section --- and the adopted flux from the
perturbation formula. On the GMAT side the same coefficients, the shared DE440
(\cref{sec:xval-constants}) and \lupnt{}'s constants are forced onto the tool, so
a residual is an integrator or frame-realisation difference only. The SRP orbit
is chosen high and sunward so the spacecraft is never eclipsed: \lupnt{}'s
apparent-disk penumbra and GMAT's shadow model differ, and a continuously sunlit
arc keeps that difference out of the cannonball comparison.

\begin{table}[htbp]
  \centering
  \caption{Force-model differences. The Orekit columns are formula-level
    acceleration checks (implemented in
    \code{cpp/test/orekit/test\_force\_models\_orekit.cc}); the GMAT columns are
    propagation-level state checks over a \SI{2}{\hour} arc (implemented in
    \code{cpp/test/gmat/test\_force\_models\_gmat.cc}); the models are specified
    in \cref{chap:dynamics}.}
  \label{tab:xval-forces}
  \scriptsize
  \begin{tabularx}{\textwidth}{@{}l W{0.62} W{0.62} W{1.76}@{}}
    \toprule
    Model & versus Orekit (accel., max) & versus GMAT (state, max) & Tolerance and origin \\
    \midrule
    Spherical-harmonic gravity &
      \SI{6.2e-15}{\meter\per\second\squared} (Earth $8\times8$ EGM96, Moon
      $12\times12$ GRGM1200B; 42 cases) &
      \SI{0.19}{\meter} over \SI{2}{\hour} (Earth $8\times8$) &
      \SI{1e-9}{\meter\per\second\squared} (Orekit); \SI{3}{\meter} /
      \SI{5e-3}{\meter\per\second} (GMAT). Orekit: matched coefficients via the
      \code{.cof}$\rightarrow$ICGEM bridge, so the residual is machine-precision
      evaluation of the same recursion. GMAT: the shared \code{.cof} and DE440
      leave only integrator and ITRF-versus-GMAT-\code{EarthFixed}
      realisation. \\
    Third body (Sun + Moon) &
      \SI{1.9e-18}{\meter\per\second\squared} (4 cases) &
      \SI{0.004}{\meter} over \SI{2}{\hour} &
      \SI{1e-12}{\meter\per\second\squared} (Orekit); \SI{0.2}{\meter} /
      \SI{1e-3}{\meter\per\second} (GMAT). The perturber positions come from the
      shared DE440 on both sides, so the point-mass formula is compared in
      isolation. \\
    Solar radiation pressure &
      \SI{6.7e-24}{\meter\per\second\squared} (cannonball, sunlit; 3 cases) &
      \SI{0.003}{\meter} over \SI{2}{\hour} (high sunlit orbit) &
      \SI{1e-15}{\meter\per\second\squared} (Orekit); \SI{0.2}{\meter} /
      \SI{1e-3}{\meter\per\second} (GMAT). Matched reference pressure and Sun
      position; the sunlit orbit keeps the two differing shadow models from
      engaging. \\
    \bottomrule
  \end{tabularx}
\end{table}

The Orekit residuals sit at machine precision --- $\num{6e-15}$,
$\num{2e-18}$ and $\num{7e-24}\,\si{\meter\per\second\squared}$ --- exactly what
a comparison of two implementations of the same closed form on identical inputs
should produce. The tolerances
(\SIlist{1e-9;1e-12;1e-15}{\meter\per\second\squared}) sit far above them yet
remain many orders of magnitude below any real force-model error. The GMAT
residuals are sub-metre over two hours because every
input is matched and only the integrator and the frame realisation differ; the
tolerances retain roughly a $15\times$ to $60\times$ margin.

\subsection{Live comparisons against SPICE}
\label{sec:xval-spice}

The SPICE comparisons are structurally different from the two above. CSPICE is
linked into \lupnt{}, so no fixture is needed: the tests call
\code{sxform\_c} or \cpp{spice::ConvertTime} and compare against \lupnt{}'s
own fitted fast paths in the same process, at the same epochs, in the same
time scale. What is being validated is not an independent implementation of a
model but the \emph{fidelity of \lupnt{}'s fits} to the SPICE kernels ---
\cpp{InitFrameConversionFromSpice} for orientation and
\cpp{InitTtMinusTdbFit} for $\TT-\TDB$ --- which exist because CSPICE is not
thread safe and cannot be called from an OpenMP region.

\begin{table}[htbp]
  \centering
  \caption{Live SPICE comparisons \citep{acton1996spice}. Probe vectors for
    the orientation checks are unit vectors, so an absolute position
    difference is numerically equal to the rotation-matrix element error.}
  \label{tab:xval-spice}
  \scriptsize
  \begin{tabularx}{\textwidth}{@{}L L l L@{}}
    \toprule
    Check & SPICE reference & Tolerance & Meaning \\
    \midrule
    \cpp{GcrfToItrf} after \cpp{InitFrameConversionFromSpice} over a 3-day window &
      \code{sxform\_c("J2000", "ITRF93", et)} &
      \num{1e-6} &
      The fitted Earth orientation reproduces the high-accuracy binary PCK to
      within \SI{1e-6}{\radian}-scale rotation error. \\
    \cpp{MoonCiToPa} after the same fit &
      \code{sxform\_c("J2000", "IAU\_MOON", et)} &
      \num{1e-6} &
      Same, for the lunar principal-axes orientation. \\
    \cpp{ItrfToGcrf}$\,\circ\,$\cpp{GcrfToItrf} and
      \cpp{MoonPaToCi}$\,\circ\,$\cpp{MoonCiToPa} &
      identity &
      \num{1e-12} &
      Round-trip closure. Angle-based fitting composes a ZXZ rotation, so the
      fitted matrix is exactly orthonormal by construction. \\
    $\mat R\transpose \mat R = \mat I$ for the fitted rotation &
      identity &
      \num{1e-14} &
      Explicit orthonormality. This would have failed at ${\sim}\num{4e-8}$
      under the earlier matrix-element fitting scheme; it is the property that
      motivated fitting angles instead. \\
    \cpp{ComputeEopFromSpice} re-composed rotation &
      \code{sxform\_c("J2000", "ITRF93", et)} &
      \num{1e-9} &
      The EOP-like parameters extracted from SPICE regenerate the original
      SPICE rotation element-wise. \\
    \cpp{spice::ConvertTime} $\TDB \leftrightarrow \TT$ round trip &
      itself &
      \SI{1e-9}{\second} &
      Round-trip closure of the DE440t path near J2000. \\
    DE440t $\TT-\TDB$ versus the truncated analytic series &
      \code{unitim\_c}-style series &
      $> \SI{1e-6}{\second}$ and $< \SI{1e-3}{\second}$ &
      The two must differ meaningfully --- that difference is the fidelity
      gain from the time ephemeris --- while both remain sane $\TT-\TDB$
      models. \\
    Native Chebyshev $\TT-\TDB$ fit (16-day segments, 13 coefficients) &
      DE440t via \cpp{spice::ConvertTime} &
      $< \SI{1e-9}{\second}$ &
      The thread-safe fast path is sub-nanosecond-faithful to the kernel. \\
    Fitted \cpp{TDBToTt} over a 2024--2025 mission window ($-30$ to $+400$ days) &
      \cpp{spice::ConvertTime} &
      $\leq 1$ ULP (${\sim}\SI{0.3}{\micro\second}$) &
      End-to-end absolute epochs agree to the last representable bit. \\
    \bottomrule
  \end{tabularx}
\end{table}

The tolerance in the last row is expressed as a ULP rather than an absolute
time because it must be: at a 2025 epoch, seconds-since-J2000 is
${\sim}\num{8e8}$, and $2^{-52}$ of that is already
${\sim}\SI{0.2}{\micro\second}$. No conversion expressed in that
representation can do better, which is the same floating-point floor that
appears in the first row of \cref{tab:xval-time}. See
\cref{chap:time-conversions} for how \lupnt{} keeps the offset arithmetic away
from the absolute epoch where precision matters.

A separate test, \code{data.eop\_latest\_iers}, exercises the live IERS
refresh described in \cref{chap:sp3-download}. It tolerates being offline, and
it explicitly restores the bundled EOP file before returning, because the
tolerances in \cref{tab:xval-frames} are calibrated against the bundled
vintage.

% ===========================================================================
\section{Why the differences are the size they are}
\label{sec:xval-causes}
% ===========================================================================

Four mechanisms account for essentially every non-trivial number in
\cref{tab:xval-time,tab:xval-frames,tab:xval-orbits}. Each is a
specification-level statement about the models involved, not a numerical
accident.

\subsection{EOP vintage}
\label{sec:xval-eop}

The GCRF$\to$ITRF rotation specified in \cref{chap:frame-conversions}
factorises into precession-nutation, Earth rotation, and polar motion. Only
the last two depend on measured data, and both come from the IERS EOP series
\citep{iers2010}. The dominant term is the Earth rotation angle, driven by
$\UTone-\UTC$.

\lupnt{}, Orekit and GMAT each ship a different snapshot of that series, taken
at a different date. The IERS revises past values as more VLBI and GNSS data
are reduced, and \emph{predicts} future ones. A tens-of-milliseconds
disagreement in $\UTone$ is therefore normal between two valid tables, and it
maps directly into an angle,
\begin{equation}
  \label{eq:xval-ut1-angle}
  \delta\theta \;=\; \omega_\oplus \, \delta(\UTone),
  \qquad \omega_\oplus \approx \SI{7.292115e-5}{\radian\per\second},
\end{equation}
so $\delta(\UTone) = \SI{0.024}{\second}$ gives
$\delta\theta \approx \SI{1.75e-6}{\radian}$. That is precisely the
$\leq \SI{1.8e-6}{\radian}$ GMST and ERA residual in \cref{tab:xval-time}: the
sidereal-angle row is not an independent check of the angle formulas at all
--- those are identical on both sides --- it is a direct read-out of the
$\UTone$ disagreement.

A rotation error is a position error proportional to radius. For a
perpendicular component,
\begin{equation}
  \label{eq:xval-rot-scaling}
  \delta r \;\approx\; \delta\theta \, \lVert\vec r\rVert .
\end{equation}
At LEO ($\lVert\vec r\rVert \approx \SI{7e6}{\meter}$) an angle of
$\num{2e-6}$ to $\num{1.4e-5}$ radians gives tens of metres --- the
\SIrange{14}{98}{\meter} and \SIrange{34}{36}{\meter} entries. At trans-lunar
distance ($\lVert\vec r\rVert \approx \SI{3.5e8}{\meter}$) the same angles
give kilometres --- the \SI{4.4}{\kilo\meter} and \SI{1.6}{\kilo\meter}
entries. This is why the tolerance is written as
$\max(\SI{100}{\meter},\, \num{2e-5}\,\lVert\vec r\rVert)$ and not as a
constant, and why the difference grows toward 2025 epochs, where one or both
tables have crossed from measurement into prediction.

\subsection{\code{MOON\_ME} versus the analytical IAU pole}
\label{sec:xval-moon}

The Moon has two distinct body-fixed frames, both specified in
\cref{chap:frame-conversions}. The \emph{principal-axes} (PA) frame is aligned
with the principal axes of the lunar inertia tensor and is what the DE lunar
libration integration actually produces. The \emph{mean-Earth/mean-rotation}
(ME) frame is the cartographic frame in which lunar products are published;
it differs from PA by a small, essentially constant rotation of order
$\SI{100}{\arcsecond}$.

A third object is often confused with both: the IAU/WGCCRE analytical pole
model, a low-order trigonometric series for the lunar right ascension,
declination and prime meridian. Orekit's body-oriented Moon frame is that
analytical model. It is designed to \emph{approximate} the ME axes, and it
does so only to about $\SI{3e-5}{\radian}$. Applying
\cref{eq:xval-rot-scaling} at a low lunar orbit radius of
${\sim}\SI{1.9e6}{\meter}$ gives ${\sim}\SI{57}{\meter}$; at ELFO apoapsis it
gives the upper end of the \SIrange{57}{200}{\meter} range. The Orekit column
is therefore not measuring an error in \lupnt{} --- it is measuring the known
accuracy of the analytical model against the kernel-based ME frame, which is
why the tolerance is written as $\num{3e-4}\,\lVert\vec r\rVert$.

GMAT gives the complementary check. Its Luna \code{BodyFixed} system is loaded
from a SPICE frame kernel and is a DE-era \emph{principal-axes} frame, so it
is compared against \lupnt{}'s \code{MOON\_PA}. With both sides now on DE440
(\cref{sec:xval-de}) the Moon-position translation agrees at the
distribution-and-interpolation level, so the residual is orientation-limited
and the two PA frames agree at the sub-arcsecond level.

The identification is confirmed by deliberately mismatching the frames.
Comparing \lupnt{}'s PA against Orekit's IAU frame, or \lupnt{}'s ME against
GMAT's Luna \code{BodyFixed}, inflates the difference by a factor of $10$ to
$40$, to \SIrange{0.6}{3.4}{\kilo\meter} --- consistent with the
${\sim}\SI{100}{\arcsecond}$ PA-to-ME rotation and inconsistent with anything
else. Between the two comparisons, both \lupnt{} lunar body-fixed frames are
independently pinned.

\subsection{ICRF versus MJ2000Eq: the frame bias}
\label{sec:xval-bias}

GCRF and EME2000 (equivalently ICRF and MJ2000Eq) differ only by the constant
\emph{frame bias} matrix, the fixed rotation of order 20 milliarcseconds
that relates the dynamical mean equator and equinox of J2000.0 to the ICRS
axes. The IERS Conventions \citep{iers2010} give it exactly, and it is
time-independent by definition.

\lupnt{} and Orekit implement that matrix. They agree to
${\sim}\num{5e-13}\,\lVert\vec r\rVert$, i.e.\ \SI{1.8e-4}{\meter} even at
trans-lunar range --- pure floating-point rounding, which is what a
comparison of two identical constant matrices should look like.

GMAT's ICRF $\leftrightarrow$ MJ2000Eq rotation differs from that matrix by
${\sim}\SI{1e-7}{\radian}$ (about 20 milliarcseconds), and the
difference is slightly epoch-dependent, which a true constant bias cannot be.
This is a GMAT-side convention difference, not an error in either tool, but it
must be budgeted for: by \cref{eq:xval-rot-scaling} it is invisible at LEO
(sub-metre) and reaches ${\sim}\SI{54}{\meter}$ at trans-lunar range. Hence
the asymmetric tolerances --- \SI{1e-3}{\meter} against Orekit but
$\max(\SI{5}{\meter},\, \num{3e-7}\,\lVert\vec r\rVert)$ against GMAT. Any
lunar analysis that ingests GMAT-produced inertial states should treat
${\sim}\SI{50}{\meter}$ as an irreducible interface bias.

\subsection{Leap-second boundaries}
\label{sec:xval-leap}

$\UTone-\UTC$ is a sawtooth: it drifts smoothly, then jumps by exactly
\SI{1}{\second} when a leap second is inserted. \lupnt{}'s EOP table stores
the daily $\UTone-\UTC$ column and interpolates it \emph{linearly across} that
jump. Within roughly a day of an insertion the interpolant therefore returns a
value that is not merely inaccurate but wrong by a large fraction of a second.
The measured errors are $+\SI{0.938}{\second}$ one hour before the
2016-12-31 leap and $+\SI{0.094}{\second}$ six hours after the 2012-06-30
leap.

The correct treatment, which Orekit uses, is to interpolate the
\emph{continuous} quantity $\UTone-\TAI$ and reconstruct $\UTone-\UTC$ from it
using the leap-second step table --- exactly the decomposition
\cref{chap:time-conversions} specifies for the constant-offset scales. The
generator flags the affected epochs with \code{near\_leap\_second: true}, and
the tests skip the $\UTone$-dependent assertions ($\UTone-\UTC$, GMST, ERA)
there. The scales that do not depend on $\UTone$ --- $\TAI$, $\TT$, $\TDB$,
$\GPST$ --- are still verified at those epochs, so leap-second handling in the
integer-offset path is not left untested.

\subsection{DE version, and the two floating-point floors}
\label{sec:xval-de}

Two further mechanisms round out the picture.

\paragraph{Ephemeris version.} \lupnt{} and Orekit both evaluate DE440
\citep{park2021de440}, and their Earth$\to$Moon and Earth$\to$Sun vectors
agree to ${\sim}\num{5e-11}$ relative --- \SI{0.33}{\meter} for the Moon and
\SI{7.2}{\meter} for the Sun, both consistent with independent
redistributions of the same Chebyshev coefficients read by independent
interpolators. GMAT ships DE405, DE421 and DE424 and still has no DE440 even in
R2026a, so it is pointed at \lupnt{}'s own \file{de440.bsp} through
\code{EphemerisSource = 'SPICE'} (\cref{sec:xval-constants}). All three tools
therefore evaluate DE440, and the former DE421-versus-DE440 lunar offset is
gone: the Moon-centred columns of \cref{tab:xval-frames} are now the same
distribution-and-interpolation residual as the Orekit column, at the
${\sim}\num{5e-11}$-relative level, rather than a DE-version translation. The
third-body and SRP force checks (\cref{tab:xval-forces}) inherit the same shared
Sun and Moon positions.

\paragraph{Floating point.} Two distinct floors appear. \lupnt{} represents
epochs as seconds since J2000 in a double; at a 2016 epoch that is
${\sim}\num{5e8}$, so one ULP is ${\sim}\SI{1e-7}{\second}$. GMAT reports
Modified Julian dates of ${\sim}30\,000$ days, and differencing two of them
amplifies the ULP further. Nothing in either tool can produce a $\TAI-\UTC$
agreement better than that, which is why the exactly-defined offsets ---
quantities that are integers or fixed constants by definition --- still show
${\sim}\SI{2e-8}{\second}$ and ${\sim}\SI{6e-7}{\second}$ residuals.

\paragraph{Integrator.} Finally, GMAT integrates numerically even where an
analytical solution exists. Its ``two-body'' propagation is RK89 with
\code{Accuracy = 1e-13}, so the \code{kepler} comparison carries genuine
integrator drift (${\sim}\SI{2e-2}{\meter}$ and ${\sim}\SI{5e-8}{\radian}$ in
the fast angles over 1.5 orbits), whereas Orekit's analytical
\code{KeplerianPropagator} agrees with \lupnt{}'s analytical propagator at
machine precision. In the inertial-axis $J_2$ case the roles reverse: both
sides integrate, and the ${\sim}\SI{2.2e-5}{\meter}$ residual is \lupnt{}'s
fixed-step RK4 at \SI{1}{\second} against Orekit's adaptive DP853 (see
\cref{chap:dynamics}).

% ===========================================================================
\section{Known modelling differences and limitations}
\label{sec:xval-limitations}
% ===========================================================================

The cross-validation surfaced three items worth stating as standing
limitations rather than as test artifacts. All three are documented in the
corresponding test files.

\paragraph{UT1 interpolation across leap seconds (a \lupnt{} limitation).}
As described in \cref{sec:xval-leap}, the daily $\UTone-\UTC$ column is
interpolated straight through the \SI{1}{\second} step, producing spurious
$\UTone-\UTC$ of up to $+\SI{0.938}{\second}$ within about a day of an
insertion. Anything that consumes $\UTone$ --- GMST, ERA, and therefore the
whole GCRF$\leftrightarrow$ITRF chain --- is affected in that window. There
have been no leap seconds since 2016-12-31, so present-day work is unaffected,
but historical reprocessing across 2012-06-30 or 2016-12-31 is not.

\paragraph{$J_2$ symmetry axis (an intentional simplification).}
\cpp{JToCartTwoBodyDynamics} applies the $J_2$ acceleration about the inertial
$Z$ axis rather than the body-fixed pole. The GMAT comparison quantifies the
cost for Earth orbits: tens of metres to ${\sim}\SI{200}{\meter}$ over up to
1.5 orbits for the tested LEO and MEO-HEO cases, growing with propagation
span. The simplification buys a frame-independent, autodiff-friendly closed
form. Use the full spherical-harmonic gravity dynamics of
\cref{chap:dynamics} when it matters.

\paragraph{GMAT's ICRF/MJ2000Eq frame bias (an external convention
difference).} As described in \cref{sec:xval-bias}, GMAT's rotation differs
from the IERS frame-bias matrix used by \lupnt{} and Orekit (which agree to
sub-millimetre) by ${\sim}\SI{1e-7}{\radian}$. It is invisible in Earth orbit
and worth ${\sim}\SI{50}{\meter}$ at trans-lunar range.

\subsection{GMAT-specific caveats}
\label{sec:xval-gmat-caveats}

Five further caveats constrain what the GMAT fixture can and cannot show.
They are recorded in \code{\_meta.notes} of
\file{cpp/test/gmat/data/gmat_reference.json} and reproduced here because
anyone extending the generator will meet them.

\begin{enumerate}[leftmargin=1.6em]
  \item \textbf{No $\UTone$ or GPS time scale.} GMAT R2026a exposes no
    \code{UT1ModJulian} or \code{GPSModJulian} spacecraft parameter, so the
    \code{time\_scales} section carries only \code{tai\_minus\_utc},
    \code{tt\_minus\_tai} and \code{tdb\_minus\_tai}.
  \item \textbf{Anomalies are Python-computed}, as described in
    \cref{sec:xval-scope}; only \code{cart0} and the propagated states are
    GMAT outputs.
  \item \textbf{DE440 through SPICE, on both sides}, with the consequences in
    \cref{sec:xval-de}: GMAT has no native DE440, so it reads \lupnt{}'s
    \file{de440.bsp}.
  \item \textbf{Luna \code{BodyFixed} is principal axes}, hence the comparison
    against \code{MOON\_PA} rather than \code{MOON\_ME}
    (\cref{sec:xval-moon}).
  \item \textbf{The ``combined \code{Propagate}'' bug.} Synchronising two
    spacecraft and propagators in one command, e.g.
    \code{Propagate Prop1(SatA) Prop2(SatB) \{SatA.ElapsedSecs = X\};}, was
    found to desynchronise the second spacecraft's force model in GMAT:
    its state ended up about $60\times$ further from the two-body solution
    than the expected $J_2$ perturbation --- clearly a tool bug, not a real
    effect. The generator therefore always emits separate, sequential
    \code{Propagate <Prop>(<Sat>) \{<Sat>.ElapsedSecs = <dt>\};} commands, one
    per spacecraft-propagator pair. Preserve that pattern when extending it.
\end{enumerate}

% ===========================================================================
\section{Regenerating the fixtures and re-running the comparison}
\label{sec:xval-regenerate}
% ===========================================================================

Regeneration is needed only when adding cross-validation cases or
deliberately updating the reference values after an intentional algorithm
change. The fixtures are checked in precisely so that nobody else has to do
it.

\begin{table}[htbp]
  \centering
  \caption{Files involved in the fixture-based comparison.}
  \label{tab:xval-files}
  \small
  \begin{tabularx}{\textwidth}{@{}L L@{}}
    \toprule
    Path & Role \\
    \midrule
    \file{cpp/test/orekit/gen_orekit_reference.py} & Generator; needs the optional \code{dev} pixi environment (JVM, JCC-wrapped Orekit). \\
    \file{cpp/test/orekit/data/orekit_reference.json} & Checked-in Orekit fixture. \\
    \file{cpp/test/orekit/test_time_conversions_orekit.cc} & Time scales, GMST, ERA. \\
    \file{cpp/test/orekit/test_frame_conversions_orekit.cc} & GCRF, EME2000, ITRF and the Moon frames. \\
    \file{cpp/test/orekit/test_ephemeris_orekit.cc} & Sun and Moon DE440 ephemerides. \\
    \file{cpp/test/orekit/test_orbit_dynamics_orekit.cc} & Kepler propagation, $J_2$ acceleration and propagation. \\
    \file{cpp/test/orekit/test_force_models_orekit.cc} & Gravity, third-body and SRP accelerations (formula level). \\
    \file{cpp/test/orekit/README.md} & Per-section data dictionary and tolerance rationale. \\
    \file{cpp/test/gmat/gen_gmat_reference.py} & Generator (time / frame / kepler / $J_2$); needs a local GMAT R2026a desktop install. \\
    \file{cpp/test/gmat/gen_gmat_forces_reference.py} & Generator (gravity, third-body, SRP propagation); same GMAT install, pointed at \lupnt{}'s DE440. \\
    \file{cpp/test/gmat/data/gmat_reference.json} & Checked-in GMAT fixture. \\
    \file{cpp/test/gmat/data/gmat_forces_reference.json} & Checked-in GMAT force-model fixture. \\
    \file{cpp/test/gmat/data/gmat_reference.script} & The exact GMAT script that was run, paths tokenised. \\
    \file{cpp/test/gmat/data/gmat_forces_reference.script} & The force-model GMAT script, paths tokenised. \\
    \file{cpp/test/gmat/lupnt_j2_earth.cof} & Minimal potential file forcing \lupnt{}'s \code{GM}, \code{R\_eq} and \code{C20} onto GMAT. \\
    \file{cpp/test/gmat/test_time_conversions_gmat.cc} & $\TAI$, $\TT$, $\TDB$. \\
    \file{cpp/test/gmat/test_frame_conversions_gmat.cc} & Earth and Moon frames. \\
    \file{cpp/test/gmat/test_orbit_dynamics_gmat.cc} & Kepler and $J_2$ propagation. \\
    \file{cpp/test/gmat/test_force_models_gmat.cc} & Gravity, third-body and SRP propagation. \\
    \file{cpp/test/gmat/README.md} & Per-section data dictionary, caveats, tolerance rationale. \\
    \file{cpp/test/utils.cc} & \cpp{LoadTestJson()}, the fixture reader. \\
    \bottomrule
  \end{tabularx}
\end{table}

\begin{lstlisting}[language=shell, caption={Regenerating the Orekit fixture.}, label={lst:xval-orekit}]
# The optional `dev` environment adds the JVM-based orekit package (via JCC)
# on top of the default environment.  Normal users and CI never need it.
pixi install -e dev
pixi run -e dev python cpp/test/orekit/gen_orekit_reference.py

# First run downloads the small `orekit-data` bundle (leap seconds, EOP, DE440)
# into data/orekit-data-main/, which .gitignore already covers via **/*orekit-data*/.
# The script overwrites cpp/test/orekit/data/orekit_reference.json.

git diff --stat cpp/test/orekit/data/orekit_reference.json   # review before committing
pixi run test-cpp
\end{lstlisting}

\begin{lstlisting}[language=shell, caption={Regenerating the GMAT fixture.}, label={lst:xval-gmat}]
# GMAT R2026a is a desktop application, not a pixi/conda package: install it
# separately.  On macOS the default location is /Applications/GMAT_R2026a/.
# Point the generator at GmatConsole via GMAT_CONSOLE, or at the install
# directory via GMAT_HOME / GMAT_ROOT_DIR; common locations are searched.
# GMAT has no native DE440, so both generators point it at LuPNT's de440.bsp
# via SolarSystem.EphemerisSource = 'SPICE'.  Standard library only.
GMAT_CONSOLE=/Applications/GMAT_R2026a/bin/GmatConsole-R2026a \
    python3 cpp/test/gmat/gen_gmat_reference.py

# The force-model fixture is a separate generator (same GMAT install; it finds
# de440.bsp under LuPNT_data by default, or set LUPNT_DE440=<path>):
GMAT_CONSOLE=/Applications/GMAT_R2026a/bin/GmatConsole-R2026a \
    python3 cpp/test/gmat/gen_gmat_forces_reference.py

# The first overwrites data/gmat_reference.json and data/gmat_reference.script;
# the second overwrites data/gmat_forces_reference.json (and .script).

pixi run test-cpp
\end{lstlisting}

Re-running the comparison alone needs neither tool:

\begin{lstlisting}[language=shell, caption={Running the cross-validation tests.}, label={lst:xval-run}]
pixi run test-cpp                      # everything, including the fixture-based tests

# Or, directly through ctest, a single family at a time:
ctest --test-dir build-cpp-test-pixi --output-on-failure -R 'orekit_reference'
ctest --test-dir build-cpp-test-pixi --output-on-failure -R 'gmat_reference'
ctest --test-dir build-cpp-test-pixi --output-on-failure -R 'conversions.frame_conversions'
ctest --test-dir build-cpp-test-pixi --output-on-failure -R 'conversions.tt_tdb_high_fidelity'
\end{lstlisting}

The relevant Catch2 test-case names are
\longid{dynamics.kepler_orekit_reference},
\longid{dynamics.j2_acceleration_orekit_reference},
\longid{dynamics.j2_propagation_orekit_reference},
\longid{dynamics.gravity_field_orekit_reference},
\longid{dynamics.third_body_orekit_reference},
\longid{dynamics.srp_orekit_reference},
\longid{dynamics.kepler_gmat_reference},
\longid{dynamics.j2_propagation_gmat_reference},
\longid{dynamics.gravity_propagation_gmat_reference},
\longid{dynamics.third_body_propagation_gmat_reference},
\longid{dynamics.srp_propagation_gmat_reference},
\longid{conversions.frame_conversions.spice_fit} and
\longid{conversions.tt_tdb_high_fidelity}.

\begin{lupntwarning}
If a regenerated fixture moves a number, do not simply widen the tolerance.
Every tolerance in
\cref{tab:xval-time,tab:xval-frames,tab:xval-orbits,tab:xval-forces} is
paired with a comment naming the mechanism that sets it. A change means either
that mechanism changed --- a newer EOP snapshot in \code{orekit-data}, a
different GMAT build, a new DE --- or that something regressed. Identify which
before touching the assertion, and update the corresponding row here.
\end{lupntwarning}

% ===========================================================================
\section{Model boundaries}
\label{sec:xval-boundaries}
% ===========================================================================

\begin{itemize}[leftmargin=1.4em]
  \item \textbf{The comparison is a snapshot.} The fixtures freeze one Orekit
    build, one \code{orekit-data} EOP vintage, and one GMAT release. They
    detect regressions in \lupnt{}; they do not track the external tools.
  \item \textbf{Coverage is by category, not exhaustive.} Three epochs, four
    Earth states and two lunar states for the frames; six orbits for
    propagation; five positions for the $J_2$ formula; 42 gravity, four
    third-body and three SRP force-model cases. Extreme regimes --- very high
    eccentricity, epochs far outside the EOP table, deep-past DE coverage ---
    are not sampled.
  \item \textbf{Drag is the one perturbation not cross-validated} against these
    tools. Point mass, $J_2$, spherical-harmonic gravity, third-body
    point-mass perturbations and solar radiation pressure are now all covered
    (\cref{tab:xval-orbits,tab:xval-forces}); atmospheric drag, eclipse/shadow
    transitions and higher-order coupling remain covered only by the internal
    tests of \cref{chap:dynamics}.
  \item \textbf{Measurement, signal and estimation models are out of scope.}
    Neither reference tool models the observables of the later chapters.
  \item \textbf{$\UTone$-dependent results are skipped near leap seconds},
    so the \SI{0.024}{\second} figure in \cref{tab:xval-time} is an
    away-from-leap-second bound, not a global one.
  \item \textbf{The SPICE comparisons validate fits, not physics.} They show
    that \lupnt{}'s thread-safe fast paths reproduce the kernels; the accuracy
    of the kernels themselves is the province of NAIF and JPL
    \citep{acton1996spice,park2021de440}.
  \item \textbf{Tolerances are calibrated against the bundled EOP vintage.}
    Refreshing EOP from the IERS mid-run (\cref{chap:sp3-download}) will move
    the ITRF comparison by tens of milliseconds of $\UTone$, i.e.\ well
    outside the asserted bound at trans-lunar range.
\end{itemize}
